\scanpage{118}
\mspage{2}
% archive sheet number: 117
% source lines 1-9
non additif, de la notion de catégorie triangulée dégagée par Verdier. Comme celle-ci, cette notion est en fait trop faible - ce n'est \uline{pas} ``la bonne'' - mais elle a un intérêt pratique provisoire. La ``bonne'' notion (s'il en est une : \unsure{p. ex|p.-ê.} celle des dérivateurs) devrait travailler avec une notion plus fine des carrés et des carrés $W$-cartésiens, en travaillant systématiquement dans
\[
\mathrm{HO}_W(\unsure{\mathbb{\Delta}|\Delta}^1\times\mathbb{\Delta}^1)=M(\mathbb{\Delta}^1\times\mathbb{\Delta}^1)W(\mathbb{\Delta}^1\times\mathbb{\Delta}^1)^{-1}
\]
\inserted{(en \unsure{y prenant|y voyant} les ``vrais'' carrés $W$-cartésiens !)} et catégories similaires, telles que $\mathrm{HO}_W(\unsure{\Psi|\Upsilon})$, $\mathrm{HO}_W(\unsure{\Phi|\varphi})$. Mais les inclure dans une axiomatique ne \unsure{faisant} pas mention d'une Cat. de modèles $(M,W)$, revient manifestement à mettre le doigt dans l'engrenage de la notion \unsure{des|de} \unsure{dérivateurs|dérivateur}, et ce n'est pas mon propos pour l'instant.

% source lines 19-24
\uline{Définition 1} J'appelle ``\uline{catégorie à carrés $h$-cartésiens}'', ou ``\uline{catégorie quadrangulée à gauche}'', une Catégorie $\unsure{\mathscr H|\mathcal H|H}$, munie d'un sous-ensemble $\mathrm{Cart}$ de l'ens $\mathrm{Hom}_{\deleted{\unsure{\mathrm{Cat}}}}(\mathbb{\Delta}^1\times\mathbb{\Delta}^1,\mathscr H)$ des diagrammes carrés
\diagram{\begin{tikzcd}[ampersand replacement=\&]
X_1 \arrow[d] \& X_3 \arrow[l] \arrow[d] \\
X_0 \& X_2 \arrow[l]
\end{tikzcd}}
dans $\mathscr H$, carrés appelés ``\uline{$h$-cartésiens}'' ($h$ comme ``homotopiquement''), Cette
